\documentclass[11pt]{article}

\usepackage[utf8]{inputenc}
\usepackage[T1]{fontenc}
\usepackage{geometry}
\geometry{margin=1in}
\usepackage{hyperref}
\hypersetup{
    colorlinks=true,
    linkcolor=blue,
    citecolor=blue,
    urlcolor=blue,
    pdftitle={AST-Based Structural Narrowing for Diff Context},
    pdfauthor={Edgar Herrador}
}
\usepackage{booktabs}
\usepackage{enumitem}
\usepackage{amsmath}
\usepackage{amssymb}
\usepackage{xcolor}
\usepackage{parskip}

\newcommand{\code}[1]{\texttt{#1}}

\title{AST-Based Structural Narrowing for Diff Context:\\
A Third Alternative to Raw Hunks and Whole Files}
\author{Edgar Herrador \\ Independent \\ \href{mailto:edgar.herrador@gmail.com}{edgar.herrador@gmail.com} \\ \href{https://github.com/eherrador}{github.com/eherrador}}
\date{}

\begin{document}

\maketitle

\begin{abstract}
Tools that derive training data or review context from code diffs --- LLM
fine-tuning pipelines built from pull-request history, automated code-review
assistants, changelog generators --- face a recurring design decision: how much
surrounding code to keep around each change. The two dominant choices, the raw
unified-diff hunk and the whole file, are both structurally mismatched to the unit
of change a developer actually reasons about: the enclosing method, function,
class, or field. We present \code{ast-context-narrowing}, a small library that
walks a Tree-sitter AST from a change's touched byte range upward to the smallest
enclosing structural definition, and document six edge-case rules --- each backed
by a regression test and, in most cases, a concrete production failure that
motivated it --- required to make that walk reliable across doc comments,
blank-line insertions, single-statement fields, and exported class declarations.
We evaluate the technique on 175 file changes drawn from 37 real merged pull
requests (103 distinct files) in a large, actively maintained open-source
TypeScript project (\code{nestjs/nest}, 76{,}705 GitHub stars at time of writing).
Narrowing resolved a definition node for every evaluated change (175/175), with a
median line-count reduction of 71.8\% relative to whole-file context --- though,
as Section~\ref{sec:results} discusses, the 100\% resolution rate reflects
TypeScript's broad definition-type set rather than evidence that narrowing is
always tight; the reduction distribution is in fact bimodal. The library, its
tests, and the evaluation harness are released under the MIT license.
\end{abstract}

\section{Introduction}

A diff-derived training or review record needs three things to be useful: the
change itself, enough surrounding code to make the change legible, and as little
\emph{irrelevant} code as possible, since every token of context is something a
downstream model (or reviewer) has to read without contributing to the signal
being conveyed. Two conventional choices for ``enough surrounding code'' both miss
this balance in opposite directions:

\begin{itemize}
\item \textbf{The raw diff hunk.} A unified diff pads the actual changed lines
with a fixed window of context (conventionally 3 lines) on each side. This window
is arbitrary relative to the code's actual structure: for a short method, it
routinely overshoots into a neighboring method or its attributes; for a large
method, it badly undershoots, omitting the method's own signature or the class it
belongs to.
\item \textbf{The whole file.} Always structurally sufficient, but often
overwhelmingly so. A one-line bug fix in a 500-line file buries the actual signal
in hundreds of unrelated lines --- costly for a human reviewer's attention, and
costly in a much more literal sense for an LLM fine-tuning pipeline, where every
token of irrelevant context dilutes the training signal and inflates the sequence
length (and therefore memory and compute cost) needed to represent each example.
\end{itemize}

We argue for a third option: recover the code's own structural unit of
change --- the method, function, class, or field the diff actually touched --- by
parsing the file and walking its AST. This is not a new idea in the abstract
(structural/AST-aware diffing has prior art, discussed in
Section~\ref{sec:related}), but we found that a naive implementation of ``find the
smallest enclosing definition node'' degrades silently back into whole-class or
whole-file context on real-world diffs unless it specifically accounts for a
handful of edge cases. Section~\ref{sec:method} documents six such rules, each
motivated by a concrete failure mode we observed narrowing a large real-world diff
corpus.

\paragraph{Contributions.}
\begin{enumerate}
\item A small, dependency-light (only Tree-sitter) library implementing
AST-based diff-to-structural-context narrowing for four languages (C\#,
TypeScript, TSX, JavaScript), released as open source.
\item Six documented, individually-tested narrowing rules, each traced to a
specific real-world failure mode.
\item An evaluation on real merged-PR diffs from a large open-source project,
quantifying both narrowing success rate and context-size reduction relative to
whole-file baseline.
\end{enumerate}

\section{Problem Statement}

Given a source file's content before a change (\code{before}), its content after
the change (\code{after}), and a unified diff between them, we want a function

\begin{quote}
\code{narrow(language, before, after, diff) -> (context, output)}
\end{quote}

such that \code{context} and \code{output} are the smallest substrings of
\code{before} and \code{after}, respectively, that (a) are syntactically
self-contained (i.e., correspond to one or more complete AST subtrees, not
arbitrary line ranges) and (b) fully contain every line the diff touched. When no
such substring exists that is smaller than the whole file --- e.g., the diff
touches only top-level statements outside any function or class, or the file fails
to parse --- the function should say so explicitly rather than silently returning
a degenerate (whole-file or empty) result, so the caller can choose its own
fallback.

\section{Method}
\label{sec:method}

\subsection{Baseline algorithm}

For each hunk in the diff, convert its touched-line range into a byte range in the
corresponding source, then walk the file's Tree-sitter AST to find the smallest
\emph{named} node whose type belongs to a per-language allowlist of ``definition''
node types (e.g., \code{method\_declaration}, \code{class\_declaration} in C\#;
\code{function\_declaration}, \code{method\_definition} in TypeScript) and that
fully contains the byte range. If a hunk's range isn't contained by any node in the
allowlist, the smallest enclosing named node of \emph{any} type is used as a
fallback signal that narrowing found no clean structural unit. The byte ranges
resolved for every hunk in a file are merged (overlapping/adjacent ranges
combined) and their source text concatenated, in source order, as the final
\code{context}/\code{output}.

This baseline is a few dozen lines of tree-walking code. It is also, on its own,
insufficient --- the following six rules were each added after the baseline
produced a visibly wrong result on a real diff.

\subsection{Rule 1 --- Anchor to touched lines, not the padded hunk span}

A unified diff hunk's line range includes a fixed window of unchanged context
lines on each side of the actual \code{+}/\code{-} lines (conventionally 3). For a
short method --- common in unit test files --- this padding alone can span past
the method's own boundaries into a neighboring method or a class-level attribute
block. When that happens, \emph{no single node} in the AST contains the padded
span, and the baseline algorithm's containment search has nothing to match at the
method level, escalating to the enclosing class.

\textbf{Fix:} re-derive the touched-line range directly from the \code{+}/\code{-}
lines in the hunk body, discarding the padding the unified-diff format adds.

\subsection{Rule 2 --- Trim trailing statement terminators}

Several grammars (JavaScript/TypeScript in particular) model a field or lexical
declaration's trailing \code{;} as a separate sibling token, not part of the
declaration node's own span. A diff line naturally includes the \code{;}
character. Left untrimmed, an otherwise-exact single-line field edit appears to
overrun the field node's boundary by one byte, which fails the ``fully contains''
containment check and escalates the match to the enclosing class for what should
be a one-line, one-node match.

\textbf{Fix:} after converting a hunk's line range to a byte range, trim a single
trailing \code{;} (and any whitespace before it) from the end of the range before
the containment search.

\subsection{Rule 3 --- Skip past a preceding doc comment}

A JSDoc block or a C\# \code{///} XML-doc summary is parsed as a sibling
\code{comment} node immediately preceding the definition it documents --- not as
part of that definition's own span. When a diff hunk rewrites both the doc comment
and the body of the method it documents in the same change (a common pattern:
updating a method's behavior alongside its documented contract), the hunk's
touched-line range starts inside the comment node. No single method-shaped node
contains a range that starts before its own node begins, so the baseline search
escalates to the enclosing class.

\textbf{Fix:} before the containment search, if the range's start byte falls
inside a comment node that is immediately followed (modulo whitespace) by more of
the touched range, advance the start byte past the comment. A hunk that touches
\emph{only} comment lines is left unmodified, since there is no following code for
the range to usefully advance into.

\subsection{Rule 4 --- A pure insertion at a blank line is a point, not a dropped hunk}

A hunk with no \code{-} lines (a pure insertion) anchors, on the ``before'' side,
to wherever the insertion occurred in the unchanged file. When that anchor line is
blank --- a common separator between two statements or methods --- naive
whitespace trimming of the byte range (needed for Rules 1--2's correctness)
collapses the range to a negative-length span, which a naive implementation drops
entirely, silently removing that hunk's contribution to the narrowed context.

\textbf{Fix:} when trimming collapses a range to zero or negative length, anchor
instead on a \emph{zero-width point} at the original (untrimmed) boundary. A point
still satisfies the containment check used by the search
(\code{node.start <= point <= node.end}), so it correctly resolves to whichever
definition the insertion actually landed inside.

\subsection{Rule 5 --- Container types (class/struct/interface) are a last resort}

A ``floating'' comment or a blank-line insertion with no enclosing
method/property/field around it (e.g., a change to a section-header comment inside
a class body, with no adjacent member touched) is, strictly speaking, contained by
the whole enclosing class --- which is \emph{technically} correct but practically
useless as narrowed context: a one-line comment edit should not pull in an entire
multi-hundred-line class.

\textbf{Fix:} partition each language's definition-type allowlist into ``member''
types (methods, functions, fields, properties) and ``container'' types (classes,
structs, interfaces). The containment search tries the smallest enclosing node of
\emph{any} type first (including unnamed/anonymous nodes), then the smallest
enclosing member type, and only falls back to a container type if no smaller
candidate of either kind exists.

\subsection{Rule 6 --- An exported container is still a container}

TypeScript/JSX register \code{export\_statement} itself as a member-level
definition type, because a bare \code{export function foo() \{\}} or
\code{export const x = ...} --- where the export wraps a non-container
declaration --- genuinely is the smallest sensible unit of context for a top-level
exported declaration. But \code{export class Foo \{ ... \}} parses as an
\code{export\_statement} node \emph{wrapping} a \code{class\_declaration}. Without
special-casing this, such a node satisfies the member-type allowlist (since
\code{export\_statement} is registered as one), sidestepping Rule 5's container
demotion entirely --- a floating comment inside an exported class would escalate
to the whole exported class instead of falling through to a smaller node.

\textbf{Fix:} at containment-check time, treat an \code{export\_statement} as a
container (subject to Rule 5's last-resort treatment) if and only if it directly
wraps a node of a container type; otherwise, treat it as a normal member.

\section{Evaluation}
\label{sec:results}

\subsection{Setup}

We ran the narrowing algorithm against every qualifying file change in the 100
most recently merged pull requests (as of evaluation time) from
\href{https://github.com/nestjs/nest}{\code{nestjs/nest}}, a large
(76{,}705-star), actively maintained, TypeScript, real-world open source project,
via the GitHub REST API. We restricted the sample to file changes where (a) the
file's extension maps to a supported language, (b) the change status is
\code{modified} (excluding wholesale file additions/deletions, which have no
meaningful ``narrow within the file'' case), (c) GitHub reported a non-empty
unified diff patch for the file, (d) the change touched fewer than 200 lines (to
keep the sample representative of typical, focused changes rather than large
mechanical refactors or generated-file updates), and (e) both file revisions were
under 60KB (to bound evaluation runtime). This yielded 175 qualifying file changes
across 103 distinct files and 37 distinct pull requests. For each we fetched the
file's full content at the PR's base and head commits and ran
\code{narrow\_diff\_context}. The evaluation script
(\code{evaluation/run\_public\_repo\_eval.py}) and raw per-file results
(\code{evaluation/results/nestjs\_nest.json}) are included in the repository for
reproduction.

\subsection{Results}

\begin{table}[h]
\centering
\begin{tabular}{lr}
\toprule
\textbf{Metric} & \textbf{Value} \\
\midrule
File changes evaluated & 175 \\
Narrowing succeeded (structural) & 175 (100.0\%) \\
No definition node found (fallback needed) & 0 \\
Parse/runtime errors & 0 \\
Mean line-count reduction (successful cases) & 53.2\% \\
Median line-count reduction (successful cases) & 71.8\% \\
Mean byte-count reduction (successful cases) & 54.0\% \\
Median byte-count reduction (successful cases) & 76.4\% \\
\bottomrule
\end{tabular}
\caption{Aggregate narrowing results on 175 real file changes from \code{nestjs/nest}.}
\label{tab:results}
\end{table}

Narrowing found \emph{some} definition node for every one of the 175 evaluated
changes. We do not read this as ``the technique reduces context by 100\% of the
time'' --- it means TypeScript's definition-type allowlist
(Section~\ref{sec:method}, \code{languages.py}) is broad enough, including
module-level \code{lexical\_declaration} and \code{export\_statement}, that almost
any top-level change lands inside \emph{some} registered definition type, even if
that definition turns out to be most of the file. The more informative signal is
the \emph{distribution} of the reduction achieved, which is bimodal rather than
tightly clustered around the median:

\begin{table}[h]
\centering
\begin{tabular}{lr}
\toprule
\textbf{Reduction} & \textbf{Share of cases} \\
\midrule
$\geq$ 90\% & 28.6\% \\
$\geq$ 75\% & 49.1\% \\
$\geq$ 50\% & 58.3\% \\
$<$ 10\% & 28.6\% \\
\bottomrule
\end{tabular}
\caption{Distribution of line-count reduction across the 175 evaluated changes.}
\label{tab:distribution}
\end{table}

Roughly 29\% of changes narrow to under a tenth of the whole file --- typically a
small, targeted edit inside one method, a single interface member, or a short
utility function in an otherwise large file (the file at the extreme end of this
group, \code{nest-application.interface.ts}, went from 211 lines to 1). At the
other end, another 29\% see under 10\% reduction --- these are overwhelmingly
cases where the diff itself is genuinely broad (e.g., a mechanical rename or
type-signature change touched consistently across a whole file, or a change to a
\code{describe} block that spans most of a test spec file), where a
whole-file-sized definition \emph{is} the correct answer, not a failure of
narrowing. The median file size in the evaluated sample was 334 lines, so even the
``no reduction'' tail is not narrowing failing silently on trivial files --- it
reflects diffs whose true scope actually is most of a substantial file.

\subsection{Threats to validity}

The evaluation covers a single project and a single (dominant) language
(TypeScript); the six rules in Section~\ref{sec:method} were each originally
motivated by C\# and TypeScript diffs from a different, private codebase during
the tool's original development (a small-language-model fine-tuning data pipeline
over a large enterprise codebase, described generically to preserve that
codebase's confidentiality), so the rules themselves are not overfit to
\code{nestjs/nest} specifically --- but the quantitative success rate and
reduction numbers above should not be assumed to transfer unchanged to C\#,
JavaScript, or TSX-heavy codebases without separate measurement. The 200-line
change-size cutoff and 60KB file-size cutoff (Section~\ref{sec:results}.1) were
chosen to keep the evaluation's runtime and API usage reasonable, not derived from
any principled threshold; a broader sample is future work.

\section{Related Work}
\label{sec:related}

Structural/AST-aware diffing has a long history --- GumTree and similar tools
compute fine-grained AST-level edit scripts between two versions of a file,
primarily for precise change \emph{detection and classification}
\cite{falleri2014finegrained} rather than context \emph{narrowing} for a
downstream consumer. Our problem is narrower and more specific: given that a
change already exists (as a unified diff, the near-universal interchange format
for VCS tooling), select the smallest well-formed structural region of the file
that contains it. This is closer in spirit to IDE ``go to enclosing symbol''
navigation than to edit-script computation, applied programmatically at data-set
or context-window construction time rather than interactively.

In LLM fine-tuning and retrieval-augmented code assistance specifically, context
construction from diffs is usually described only at the level of ``we extract the
relevant function/class'' without publishing the edge-case handling that makes such
extraction reliable across real-world diffs --- the six rules in
Section~\ref{sec:method} are, to our knowledge, not documented elsewhere as a
checklist, despite (based on our own experience discovering them one production
failure at a time) being generic to any Tree-sitter-based narrowing implementation
rather than specific to our original use case. The parsing itself is built on
Tree-sitter \cite{treesitter}, an incremental parsing library providing concrete
syntax trees for many languages.

\section{Limitations}

\begin{itemize}
\item \textbf{No hybrid-grammar support.} Razor (\code{.cshtml}), Vue single-file
components, and similar host/guest-language files are not supported; narrowing
them needs dedicated handling of the embedding boundary, which we leave to future
work.
\item \textbf{Definition-type allowlists are manually curated per language.}
Extending support to a new language requires identifying the right Tree-sitter
node types by hand; we have not attempted to derive this automatically from a
grammar.
\item \textbf{Single-file scope.} The algorithm narrows within one file at a time
and does not attempt cross-file context (e.g., pulling in a called function's
definition from another file), which some use cases may still need on top of this
technique.
\item \textbf{Language-, not natural-language-, specific.} The technique's only
dependency on the underlying grammar is syntactic; it makes no natural-language
assumptions about identifiers or comments.
\end{itemize}

\section{Conclusion}

Whole-file and raw-diff-hunk context are both structurally mismatched to how
developers actually reason about a change. AST-based narrowing recovers the code's
own unit of change, but only reliably once a handful of specific,
individually-motivated edge cases are handled. We release
\code{ast-context-narrowing}, implementing and testing all six rules, under the
MIT license, along with the evaluation harness used in Section~\ref{sec:results},
so the technique and its evaluation are both reproducible and extensible to new
languages and codebases.

\section*{Availability}

Code: \url{https://github.com/eherrador/ast-context-narrowing} (MIT license)

\bibliographystyle{plain}
\bibliography{references}

\end{document}
